% AI drafted, human edited.

% 00-summary.tex -- the results, stated without their derivations.
% Formatted as an appendix for now so that the equation numbering of the body does not
% shift while it is being drafted; it is intended to become section 0, before section 1.
% Its displays carry labels from the start, so only the counter changes when it moves.

\appendix
% article's \appendix renumbers sections but not equations, so the counter is reset here and
% given an A prefix.  When this becomes section 0 the two lines go and the numbers become
% ordinary ones; the labels are unchanged either way.
\setcounter{equation}{0}
\renewcommand{\theequation}{A\arabic{equation}}
\section{Not Really an Appendix}
\label{sec:summary}

This is the whole calculation, stated and not derived.  A reader who stops here has the
essentials; the sections that follow are the derivations, and each part of this appendix
corresponds to one of them.

Conventions throughout: the sky axes are $x$ towards North and $y$ towards East, the third
axis points at the observer, and a map is normalised so that its monopole is $a_{00} = 1$,
which is the sphere mean itself.  The star has unit angular radius, so the
visible disk is the unit circle.

\subsection{Visibility, and the expansion of a map}

The complex visibility is the Fourier transform of the sky brightness with $2\pi$ inside the
exponent and no prefactor,
\begin{equation}
  V(u,v) \;=\; \int I(x,y)\, e^{-2\pi i (u x + v y)}\, dx\, dy .
  \label{eq:summary-visibility}
\end{equation}
A real map is expanded in complex harmonics using non-negative orders only, the reality
being carried by the expansion itself,
\begin{equation}
  I(\theta,\varphi) = \sum_{l} \sum_{m=-l}^{l} a_{lm}Y_{lm}
  \label{eq:summary-expansion}
\end{equation}
The same pairing works through the transform, and the parity of $m$ decides which part of a
coefficient survives: the contribution of one stored $a_{lm}$ is twice the real part of
$a_{lm}V_{lm}$ when $m$ is even and twice $i$ times its imaginary part when $m$ is odd.  That
is the whole of the negative-order bookkeeping, and it is why half the coefficients suffice.
Its one consequence worth stating early: at a negative order the visibility carries no extra
sign, because the $(-1)^m$ of the basis is cancelled by the $(-1)^m$ that $e^{im\pi}$
contributes when the baseline is reversed.

\subsection{Rotation}

A map carried in the body frame reaches the observer through a rotation of its coefficients,
$b = D(R)a$, where $D$ is the Wigner matrix of the rotation $R$.  In the zyz convention
$R = R_z(\gamma)R_y(\beta)R_z(\alpha)$, with $\alpha$ the spin phase, $\beta$ the inclination
and $\gamma$ the sky position angle of the spin.  Two features of $D$ do the work.  It is
block diagonal in $l$, so the degree is untouched and a truncated map stays truncated.  And
each block is a finite matrix of $(2l+1)^2$ numbers that depends only on the three angles,
so a rotation costs a matrix product rather than a resampling.

For the zonal $l = 2$ patterns that the deformation needs, this specialises to one identity
and five numbers.  With $\hat e$ the axis of the pattern,
\begin{equation}
  P_2(\hat n\cdot\hat e) \;=\; \frac{1}{5}\Bigl[\, \overline{Y_{20}(\hat e)}\,Y_{20}(\hat n)
  \;+\; 2\sum_{m=1}^{2}\mathrm{Re}\bigl(\overline{Y_{2m}(\hat e)}\,Y_{2m}(\hat n)\bigr)\Bigr] ,
  \label{eq:summary-zonal}
\end{equation}
and the five numbers, for an axis at direction cosines $(a,b,c)$,
\begin{equation}
  \begin{aligned}
    Y_{20}(\hat e) &= \sqrt{\tfrac54}\,\bigl(3c^2-1\bigr), \\
    Y_{2,\pm1}(\hat e) &= \mp\sqrt{\tfrac{15}{2}}\,c\,\bigl(a\pm ib\bigr), \\
    Y_{2,\pm2}(\hat e) &= \sqrt{\tfrac{15}{8}}\,\bigl(a\pm ib\bigr)^2 .
  \end{aligned}
  \label{eq:summary-five-numbers}
\end{equation}
No zonal pattern needs more than these, whatever frame its axis is quoted in.

\subsection{The reduction to one radial integral}

Because the disk is a circle, the angular part of \eqref{eq:summary-visibility} can be done
once and for all: each harmonic contributes a complex constant times a radial integral that
depends on the baseline length $q$ alone,
\begin{equation}
  V_{lm}(u,v) \;=\; 2\pi(-i)^m N_{lm}\, e^{im\psi}\, X_{lm}(q) ,
  \label{eq:summary-mode}
\end{equation}
with $\psi$ the baseline's position angle, $N_{lm}$ the harmonic normalisation, and
\begin{equation}
  X_{lm}(q) \;=\; \int_0^1 P_{lm}\bigl(\sqrt{1-\rho^2}\bigr)\, J_m(2\pi q\rho)\,
  \rho\, d\rho .
  \label{eq:summary-radial-integral}
\end{equation}
Everything about the geometry of a single mode is in $X_{lm}$; everything about its place on
the sky is in $e^{im\psi}$.  The reduction is exact at every $q$, and it is the object the
next two subsections evaluate in closed form.

\subsection{The closed forms}

The integral splits by the parity of $l+m$, and the two halves are different in kind.

For $l+m$ odd, $P_{lm}(\sqrt{1-\rho^2})$ is $\rho^m$ times a polynomial in $\rho^2$ of degree
$(l-m-1)/2$, and the Bessel integral collapses to a single term:
\begin{equation}
  X_{lm}(q) \;=\; (-1)^m\,\frac{(l+m)!!}{(l-m-1)!!}\,
  \frac{j_l(2\pi q)}{2\pi q} , \qquad l+m \text{ odd} .
  \label{eq:summary-odd}
\end{equation}
One spherical Bessel function, whatever the degree.

For $l+m$ even there is an extra factor $\sqrt{1-\rho^2}$, and the same polynomial is instead
expanded in Zernike radial polynomials of the same order,
\begin{equation}
  \rho^m A(\rho^2) \;=\; \sum_{k=0}^{(l-m)/2} c_k\, R_{m+2k}^{m}(\rho) ,
  \qquad
  c_k \;=\; \bigl(2(m+2k)+2\bigr)\int_0^1 \rho^m A(\rho^2)\,R_{m+2k}^{m}(\rho)\,
  \rho\, d\rho .
  \label{eq:summary-even-coefficients}
\end{equation}
The expansion is finite because both sides have the same degree and parity, and the Zernike
transform turns it term by term into ordinary Bessel functions,
\begin{equation}
  X_{lm}(q) \;=\; \sum_{k=0}^{(l-m)/2} c_k\,(-1)^k\,
  \frac{J_{m+2k+1}(2\pi q)}{2\pi q} , \qquad l+m \text{ even} .
  \label{eq:summary-even}
\end{equation}
So every mode of every degree is a finite Bessel sum, and a visibility at a single baseline
is a short sum over modes with no quadrature anywhere.

The expressions \eqref{eq:summary-odd} and \eqref{eq:summary-even}
were derived by \cite{dholakia2025} but the
expression \eqref{eq:summary-even-coefficients} giving the
coefficients $c_k$ as rational numbers is a refinement.

\subsection{The deformed source}

A star deformed by its own rotation and by a companion has a level
surface given, to quadrupole order, by a potential of two amplitudes,
\begin{equation}
  \Phi(\mathbf{x}) \;=\; -\frac{1}{r}
  \;-\;\epsilon_\text{tide}\,r^2P_2(n\cdot\hat e_\text{tide})
  \;+\;\epsilon_\text{spin}\,r^2P_2(n\cdot\hat e_\text{spin})
  \;-\;\epsilon_\text{spin}\,r^2 ,
  \label{eq:summary-potential}
\end{equation}
whose surface is the level set $\Phi = -1$.  Writing the depth of the perturbation along a
direction as $E(n) = \epsilon_\text{tide}P_2(n\cdot\hat e_\text{tide})
- \epsilon_\text{spin}[P_2(n\cdot\hat e_\text{spin})-1]$, that surface is the solution of
$1/r + E(n)r^2 = 1$, which is exact and costs one root per sampled direction.

The brightness follows the gravity, $I \propto g^p$ with $g = |\nabla\Phi|$ on the surface and
$p = 1$ the bolometric von Zeipel law, multiplied by a linear limb factor
\begin{equation}
  1 - u_\text{limb}(1-\mu) , \qquad \mu = \hat n\cdot\hat e_\text{obs} ,
  \label{eq:summary-limb-law}
\end{equation}
which is linear in $\mu$ because $\mu$ is a single harmonic, and therefore costs one
coefficient at $m = 0$ rather than a product of expansions.

The shape and the brightness together make the map a function of position that no closed form
covers, so the scale of the calculation is set by how the surface is approximated.  Two
approximations are used.  The first replaces the sky map by its linear part: the shape is
captured by a single $2\times2$ positive-definite matrix, the stretch
\begin{equation}
  \Lambda \;=\; \sqrt{A A^{T}} , \qquad A = P\,(\text{frame}\cdot\text{radii}) ,
  \label{eq:summary-stretch}
\end{equation}
where $P$ projects onto $(\text{East},\text{North})$ and the radii are taken along the
principal directions.  Those directions are the eigenvectors of the quadratic form of
$E(n)$, and the radii along them are the exact level radii, so $\Lambda$ is computed from the
level surface rather than fitted.  The stretch acts on the visibility as a change of
frequency,
\begin{equation}
  V_\text{affine}(u,v) \;=\; V_\text{sph}\bigl(\Lambda^{T}(u,v)\bigr) ,
  \label{eq:summary-affine}
\end{equation}
with no flux factor: the determinant that stretching puts on the brightness cancels the one
it puts on the measure.  $\Lambda$ is best read as a preconditioner, since a different but
equally defensible stretch changes the coefficients of any series built on it without
changing the exact answer.

To first order in both amplitudes the pattern is a constant plus two quadrupoles, and with
the limb factor carried through it becomes
\begin{equation}
  I(\hat n) \;=\; \bigl(1 - u_\text{limb}(1-\mu)\bigr)\Bigl[\,1 - 4\epsilon_\text{spin}
  - 4\epsilon_\text{tide}P_2(\hat n\cdot\hat t)
  + 4\epsilon_\text{spin}P_2(\hat n\cdot\hat s)\Bigr] ,
  \label{eq:summary-first-order}
\end{equation}
whose coefficients are, to leading order in $u_\text{limb}\epsilon$,
\begin{equation}
  \begin{aligned}
    a_{00} &= (1-u_\text{limb})\,(1-4\epsilon_\text{spin}) , \\
    a_{10} &= \frac{u_\text{limb}}{\sqrt3} , \qquad
    a_{2m} = -\frac{4}{5}\Bigl[\epsilon_\text{tide}\,\overline{Y_{2m}(\hat t)}
    - \epsilon_\text{spin}\,\overline{Y_{2m}(\hat s)}\Bigr] .
  \end{aligned}
  \label{eq:summary-first-order-coefficients}
\end{equation}
A constant, a dipole from the limb, and two quadrupoles from the deformation, with no Gaunt
integral anywhere: the limb is a dipole because $\mu$ is a single harmonic, and the
quadrupoles come from \eqref{eq:summary-zonal} and \eqref{eq:summary-five-numbers}.

The second approximation does not assume the coefficients at all.  The surface is rendered,
sampled on the level surface with the limb factor applied, and coefficients fitted by least
squares in the image plane, which is the norm the data live in.  The fit solves a normal
equation whose matrix is a Gram matrix of the basis alone, and whose right-hand side is a sum
over the samples,
\begin{equation}
  b_{lm} \;=\; \sum_s w_s\, I_s\, \overline{\Phi_{lm}\bigl(\Lambda^{-1}x_s\bigr)} ,
  \label{eq:summary-fitted-rhs}
\end{equation}
with $\Phi_{lm}$ the projected harmonics and $w_s$ the projected area each sample carries.
This row has no assumption beyond the stretch, so its error is a truncation at the chosen
degree, which the user controls; the first-order row's error is instead the linearisation of
the brightness.  Comparing the two is what the interactive models are for.



\subsection{What is new here}

Sections 1--3 are standard spherical-harmonic material, and the reduction of
\eqref{eq:summary-mode}--\eqref{eq:summary-radial-integral} follows published work.  The
closed forms of \eqref{eq:summary-odd} and \eqref{eq:summary-even} use the same parity split
as Dholakia \& Pope (2025), but give each mode as an explicit Bessel sum rather than as a
basis-change matrix.  The deformation of
\eqref{eq:summary-potential}--\eqref{eq:summary-first-order-coefficients}, with the limb
carried in the coefficients, and the fitted row of \eqref{eq:summary-fitted-rhs}, are the new
material, and the two interactive models that compare them are the second.
